\subsection{REST API Details}

All error responses are serialized through \texttt{Message.toJSON()} and therefore use the envelope \texttt{\{"message":\{"message":...,"error\_code":...,"error\_details":...\}\}}; \texttt{error\_code} and \texttt{error\_details} are omitted when null. \texttt{Accept} is mandatory for every documented endpoint and must include \texttt{application/json} or be exactly \texttt{*/*}. For \texttt{POST} and \texttt{PUT}, \texttt{Content-Type} is also mandatory and must contain \texttt{application/json}; this includes \texttt{POST /user/logout} and \texttt{PUT /order/\{order\_id\}/\{order\_status\}}, even though those handlers do not read a request body.

Common globally enforced REST errors are:

\begin{itemize}
\item \texttt{400 Bad Request}, \texttt{E4V1} - inner message \texttt{Output media type not specified.}, inner details \texttt{Accept request header missing.}
\item \texttt{406 Not Acceptable}, \texttt{E4V2} - inner message \texttt{Unsupported output media type. Resources are represented only in application/json.}, inner details \texttt{Requested representation is <Accept>.}
\item \texttt{400 Bad Request}, \texttt{E4V3} - inner message \texttt{Input media type not specified.}, inner details \texttt{Content-Type request header missing.}
\item \texttt{415 Unsupported Media Type}, \texttt{E4V4} - inner message \texttt{Unsupported input media type. Resources are represented only in application/json.}, inner details \texttt{Submitted representation is <Content-Type>.}
\item Protected-route dispatcher errors:

\texttt{401 Unauthorized}, \texttt{E4A2} - inner message \texttt{Authentication required.} \texttt{403 Forbidden}, \texttt{E4A4} - inner message \texttt{You do not have permission to perform this operation.}
\item Unmatched-route dispatcher error:

\texttt{404 Not Found}, \texttt{E4R1} - inner message \texttt{Unknown resource requested.}, inner details \texttt{Requested resource is <request URI>.}
\item Generic fallback:

\texttt{500 Internal Server Error}, \texttt{E5G1} - either \texttt{Unexpected error.} from \texttt{RestDispatcherServlet} or \texttt{Unable to serve the REST request: <ACTION>.} from \texttt{AbstractRR}, with dynamic \texttt{error\_details}.
\end{itemize}

\subsubsection*{Register User}

Creates a new user account. The DAO always persists role \texttt{CUSTOMER}.

\begin{itemize}[nosep]
\item \textbf{URL:} \texttt{/user/register}

\item \textbf{Method:} \texttt{POST}

\item \textbf{URL Parameters:} \texttt{None}

\item \textbf{Query Parameters:} \texttt{None}

\item \textbf{Headers:} \texttt{Content-Type: application/json}; \texttt{Accept: application/json}

\item \textbf{Cookies:} \texttt{None}

\item \textbf{Data Parameters:}

\begin{lstlisting}
{
  "username": "newcustomer",
  "email": "newcustomer@example.com",
  "password": "Secret1!",
  "name": "Giulia",
  "surname": "Rossi",
  "phone_number": "3331234567"
}
\end{lstlisting}

Required fields are \texttt{username}, \texttt{email}, \texttt{password}, \texttt{name}, \texttt{surname}, and \texttt{phone\_number}; blank values are rejected for all of them. \texttt{email} must match \texttt{\textasciicircum{}[A-Za-z0-9+\_.-]+@[A-Za-z0-9.-]+\textbackslash{}\textbackslash{}.[A-Za-z]\{2,\}\$}. \texttt{password} must be 8-16 characters long and contain at least one lowercase letter, one uppercase letter, one digit, and one special character. The shared \texttt{User.fromJSON()} parser also accepts \texttt{id} and \texttt{role}, but this endpoint ignores them on success: the database generates the id and \texttt{RegisterUserDAO} always inserts \texttt{CUSTOMER}. An invalid \texttt{role} enum value is not explicitly handled and falls through to \texttt{E5G1}.

\item \textbf{Success Response:} \texttt{201 Created}

Content:

\begin{lstlisting}
{
  "id": 6
}
\end{lstlisting}

\item \textbf{Error Responses:}

\begin{itemize}[nosep]
\item Common header errors: \texttt{E4V1}, \texttt{E4V2}, \texttt{E4V3}, \texttt{E4V4}.
\item \texttt{400 Bad Request}, \texttt{E4V8} - inner message \texttt{Missing required fields.}, inner details \texttt{Fields 'username', 'email', 'password', 'name', 'surname', and 'phone\_number' are required.}
\item \texttt{400 Bad Request}, \texttt{E4V8} - inner message \texttt{Invalid email address.}, inner details \texttt{The provided email address is not valid.}
\item \texttt{400 Bad Request}, \texttt{E4V8} - inner message \texttt{Invalid password.}, inner details \texttt{Password must be 8-16 characters and contain at least one uppercase letter, one lowercase letter, one digit, and one special character.}
\item \texttt{400 Bad Request}, \texttt{E4V8} - inner message \texttt{Malformed request body.}, inner details are the parser message or \texttt{Unexpected key in JSON: <field>}.
\item \texttt{409 Conflict}, \texttt{E4R3} - inner message \texttt{Email, username, or phone number already in use.}
\item \texttt{500 Internal Server Error}, \texttt{E5D1} - inner message \texttt{Registration failed: database error.}, inner details are the SQL exception message.
\item \texttt{500 Internal Server Error}, \texttt{E5G1} - uncaught runtime failures, including invalid \texttt{role} enum values accepted by the shared parser.
\end{itemize}

\item \textbf{Sample Request Body:}

\begin{lstlisting}
{
  "username": "newcustomer",
  "email": "newcustomer@example.com",
  "password": "Secret1!",
  "name": "Giulia",
  "surname": "Rossi",
  "phone_number": "3331234567"
}
\end{lstlisting}
\end{itemize}

\subsubsection*{Authenticate User}

Authenticates an existing user and issues a signed JWT in the \texttt{auth\_token} HttpOnly cookie.

\begin{itemize}[nosep]
\item \textbf{URL:} \texttt{/user/login}

\item \textbf{Method:} \texttt{POST}

\item \textbf{URL Parameters:} \texttt{None}

\item \textbf{Query Parameters:} \texttt{None}

\item \textbf{Headers:} \texttt{Content-Type: application/json}; \texttt{Accept: application/json}

\item \textbf{Cookies:} \texttt{None}

\item \textbf{Data Parameters:}

\begin{lstlisting}
{
  "email": "merja@test.com",
  "password": "@UniPadova26"
}
\end{lstlisting}

Only \texttt{email} and \texttt{password} are parsed. No controller-level e-mail format validation is performed here; the credentials are passed directly to \texttt{AuthenticateUserDAO}. On success, the JWT contains the claims \texttt{user\_id} and \texttt{user\_role}, expires after 8 hours (\texttt{28800} seconds), and is sent as \texttt{Set-Cookie: auth\_token=<JWT>; Path=/; HttpOnly; SameSite=Strict; Max-Age=28800}. No \texttt{Secure} attribute is added in code.

\item \textbf{Success Response:} \texttt{200 OK}

Content: \texttt{None} Additional header: \texttt{Set-Cookie: auth\_token=<JWT>; Path=/; HttpOnly; SameSite=Strict; Max-Age=28800}

\item \textbf{Error Responses:}

\begin{itemize}[nosep]
\item Common header errors: \texttt{E4V1}, \texttt{E4V2}, \texttt{E4V3}, \texttt{E4V4}.
\item \texttt{400 Bad Request}, \texttt{E4V8} - inner message \texttt{Missing email or password.}, inner details \texttt{Fields 'email' and 'password' are required.}
\item \texttt{400 Bad Request}, \texttt{E4V8} - inner message \texttt{Malformed request body.}, inner details are \texttt{Unexpected key in JSON: <field>}.
\item \texttt{400 Bad Request}, \texttt{E4V6} - inner message \texttt{Malformed JSON in request body.}, inner details are the Jackson parser message.
\item \texttt{401 Unauthorized}, \texttt{E4A1} - inner message \texttt{Invalid email or password.}
\item \texttt{500 Internal Server Error}, \texttt{E5D1} - inner message \texttt{Login failed: database error.}, inner details are the SQL exception message.
\end{itemize}

\item \textbf{Sample Request Body:}

\begin{lstlisting}
{
  "email": "merja@test.com",
  "password": "@UniPadova26"
}
\end{lstlisting}
\end{itemize}

\subsubsection*{Logout User}

Expires the JWT cookie.

\begin{itemize}[nosep]
\item \textbf{URL:} \texttt{/user/logout}

\item \textbf{Method:} \texttt{POST}

\item \textbf{URL Parameters:} \texttt{None}

\item \textbf{Query Parameters:} \texttt{None}

\item \textbf{Headers:} \texttt{Content-Type: application/json; Accept: application/json}

\item \textbf{Cookies:} \texttt{auth\_token} HttpOnly cookie containing a valid JWT

\item \textbf{Data Parameters:} \texttt{None}

Although the handler does not read a body, the shared \texttt{AbstractRR} logic still requires \texttt{Content-Type: application/json} because the route uses \texttt{POST}.

\item \textbf{Success Response:} \texttt{204 No Content}

Content: \texttt{None} Additional header: \texttt{Set-Cookie: auth\_token=; Path=/; HttpOnly; SameSite=Strict; Max-Age=0}

\item \textbf{Error Responses:}

\begin{itemize}[nosep]
\item Common auth error: \texttt{401 Unauthorized}, \texttt{E4A2}. Invalid or expired JWTs are not mapped to \texttt{E4A3}; they are treated as unauthenticated and produce \texttt{E4A2}.
\item Common header errors: \texttt{E4V1}, \texttt{E4V2}, \texttt{E4V3}, \texttt{E4V4}.
\end{itemize}

\end{itemize}

\subsubsection*{Get Current User}

Returns the profile of the authenticated user identified by the JWT \texttt{user\_id} claim.

\begin{itemize}[nosep]
\item \textbf{URL:} \texttt{/user}

\item \textbf{Method:} \texttt{GET}

\item \textbf{URL Parameters:} \texttt{None}

\item \textbf{Query Parameters:} \texttt{None}

\item \textbf{Headers:} \texttt{Accept: application/json}

\item \textbf{Cookies:} \texttt{auth\_token} HttpOnly cookie containing a valid JWT

\item \textbf{Data Parameters:} \texttt{None}

The endpoint does not accept a path id. It always resolves the current user from the dispatcher-set \texttt{user\_id} request attribute.

\item \textbf{Success Response:} \texttt{200 OK}

Content:

\begin{lstlisting}
{
  "id": 4,
  "username": "gpadoan",
  "email": "padoan@test.com",
  "name": "Giancarlo",
  "surname": "Padoan",
  "phone_number": "1234567",
  "role": "CUSTOMER"
}
\end{lstlisting}

\item \textbf{Error Responses:}

\begin{itemize}[nosep]
\item Common auth error: \texttt{401 Unauthorized}, \texttt{E4A2}.
\item Common header errors: \texttt{E4V1}, \texttt{E4V2}.
\item \texttt{404 Not Found}, \texttt{E4R2} - inner message \texttt{User with id 4 not found.} when the JWT user id no longer exists in the database.
\item \texttt{400 Bad Request}, \texttt{E4V6} - inner message \texttt{Invalid user ID: <value>} if the dispatcher-set \texttt{user\_id} attribute is not numeric.
\item \texttt{500 Internal Server Error}, \texttt{E5D1} - inner message \texttt{Database error while retrieving user.}, inner details are the SQL exception message.
\end{itemize}

\end{itemize}

\subsubsection*{Edit Current User}

Updates the authenticated user's profile.

\begin{itemize}[nosep]
\item \textbf{URL:} \texttt{/user}

\item \textbf{Method:} \texttt{PUT}

\item \textbf{URL Parameters:} \texttt{None}

\item \textbf{Query Parameters:} \texttt{None}

\item \textbf{Headers:} \texttt{Content-Type: application/json}; \texttt{Accept: application/json}

\item \textbf{Cookies:} \texttt{auth\_token} HttpOnly cookie containing a valid JWT

\item \textbf{Data Parameters:}

\begin{lstlisting}
{
  "name": "Giancarlo",
  "surname": "Padoan",
  "email": "giancarlo.padoan@example.com",
  "phone_number": "3339876543"
}
\end{lstlisting}

The endpoint uses the JWT \texttt{user\_id}; it does not accept a target user id in the URL. Only \texttt{name}, \texttt{surname}, \texttt{email}, and \texttt{phone\_number} are actually required by the handler, and it checks only for \texttt{null}, not blank strings. No controller-level e-mail-format validation is performed here. The shared \texttt{User.fromJSON()} parser also accepts \texttt{id}, \texttt{username}, \texttt{role}, and \texttt{password}, but this handler ignores them on success; invalid \texttt{role} enum values and malformed JSON syntax are not explicitly caught and fall through to \texttt{E5G1}.

\item \textbf{Success Response:} \texttt{200 OK}

Content:

\begin{lstlisting}
{
  "id": 4,
  "username": "gpadoan",
  "email": "giancarlo.padoan@example.com",
  "name": "Giancarlo",
  "surname": "Padoan",
  "phone_number": "3339876543",
  "role": "CUSTOMER"
}
\end{lstlisting}

\item \textbf{Error Responses:}

\begin{itemize}[nosep]
\item Common auth error: \texttt{401 Unauthorized}, \texttt{E4A2}.
\item Common header errors: \texttt{E4V1}, \texttt{E4V2}, \texttt{E4V3}, \texttt{E4V4}.
\item \texttt{400 Bad Request}, \texttt{E4V6} - inner message \texttt{Missing required fields: name, surname, email and phone\_number must be provided.}
\item \texttt{400 Bad Request}, \texttt{E4V8} - inner message \texttt{Malformed JSON in request body.}, inner details are \texttt{Unexpected key in JSON: <field>}.
\item \texttt{404 Not Found}, \texttt{E4R2} - inner message \texttt{User with id 4 not found.}
\item \texttt{409 Conflict}, \texttt{E4R3} - inner message \texttt{Email or phone number already in use.}
\item \texttt{500 Internal Server Error}, \texttt{E5D1} - inner message \texttt{Database error while updating user.}, inner details are the SQL exception message.
\item \texttt{500 Internal Server Error}, \texttt{E5G1} - malformed JSON syntax and invalid \texttt{role} enum values fall through to \texttt{Unable to serve the REST request: EDIT\_USER.}
\end{itemize}

\item \textbf{Sample Request Body:}

\begin{lstlisting}
{
  "name": "Giancarlo",
  "surname": "Padoan",
  "email": "giancarlo.padoan@example.com",
  "phone_number": "3339876543"
}
\end{lstlisting}
\end{itemize}

\subsubsection*{Delete Current User}

Deletes the authenticated user's account.

\begin{itemize}[nosep]
\item \textbf{URL:} \texttt{/user}

\item \textbf{Method:} \texttt{DELETE}

\item \textbf{URL Parameters:} \texttt{None}

\item \textbf{Query Parameters:} \texttt{None}

\item \textbf{Headers:} \texttt{Accept: application/json}

\item \textbf{Cookies:} \texttt{auth\_token} HttpOnly cookie containing a valid JWT

\item \textbf{Data Parameters:} \texttt{None}

The endpoint uses the JWT \texttt{user\_id}. Because the schema defines \texttt{orders.user\_id} with \texttt{ON DELETE CASCADE}, deleting a user also deletes that user's orders. The handler does not clear the auth cookie after deletion.

\item \textbf{Success Response:} \texttt{200 OK}

Content:

\begin{lstlisting}
{
  "id": 4,
  "username": "gpadoan",
  "email": "padoan@test.com",
  "name": "Giancarlo",
  "surname": "Padoan",
  "phone_number": "1234567",
  "role": "CUSTOMER"
}
\end{lstlisting}

\item \textbf{Error Responses:}

\begin{itemize}[nosep]
\item Common auth error: \texttt{401 Unauthorized}, \texttt{E4A2}.
\item Common header errors: \texttt{E4V1}, \texttt{E4V2}.
\item \texttt{404 Not Found}, \texttt{E4R2} - inner message \texttt{User with id 4 not found.}
\item \texttt{400 Bad Request}, \texttt{E4V6} - inner message \texttt{Invalid user ID: <value>} if the dispatcher-set \texttt{user\_id} is not numeric.
\item \texttt{500 Internal Server Error}, \texttt{E5D1} - inner message \texttt{Database error while deleting user.}, inner details are the SQL exception message.
\end{itemize}

\end{itemize}

\subsubsection*{Create Category}

Creates a new category.

\begin{itemize}[nosep]
\item \textbf{URL:} \texttt{/category}

\item \textbf{Method:} \texttt{POST}

\item \textbf{URL Parameters:} \texttt{None}

\item \textbf{Query Parameters:} \texttt{None}

\item \textbf{Headers:} \texttt{Content-Type: application/json}; \texttt{Accept: application/json}

\item \textbf{Cookies:} \texttt{auth\_token} HttpOnly cookie containing a valid JWT with role \texttt{ADMIN}

\item \textbf{Data Parameters:}

\begin{lstlisting}
{
  "name": "Pizze"
}
\end{lstlisting}

Only the \texttt{name} field is parsed. \texttt{name} is required and blank values are rejected. Unexpected fields are mapped to \texttt{E4V8}. Malformed JSON syntax is not explicitly caught here and therefore falls through to \texttt{E5G1}.

\item \textbf{Success Response:} \texttt{201 Created}

Content:

\begin{lstlisting}
{
  "name": "Pizze"
}
\end{lstlisting}

\item \textbf{Error Responses:}

\begin{itemize}[nosep]
\item Common auth errors: \texttt{401 Unauthorized}, \texttt{E4A2}; \texttt{403 Forbidden}, \texttt{E4A4}.
\item Common header errors: \texttt{E4V1}, \texttt{E4V2}, \texttt{E4V3}, \texttt{E4V4}.
\item \texttt{400 Bad Request}, \texttt{E4V6} - inner message \texttt{Missing required field: name must be provided.}
\item \texttt{400 Bad Request}, \texttt{E4V8} - inner message \texttt{Malformed JSON in request body.}, inner details are \texttt{Unexpected field: <field>}.
\item \texttt{409 Conflict}, \texttt{E5D1} - inner message \texttt{Category already exists.}, inner details are the database exception message. The status is \texttt{409}, but the code is the database code \texttt{E5D1}, not \texttt{E4R3}.
\item \texttt{500 Internal Server Error}, \texttt{E5D1} - inner message \texttt{Database error while creating category.}, inner details are the SQL exception message.
\item \texttt{500 Internal Server Error}, \texttt{E5G1} - malformed JSON syntax falls through to \texttt{Unable to serve the REST request: CREATE\_CATEGORY.}
\end{itemize}

\item \textbf{Sample Request Body:}

\begin{lstlisting}
{
  "name": "Pizze"
}
\end{lstlisting}
\end{itemize}

\subsubsection*{Get All Categories}

Returns all categories.

\begin{itemize}[nosep]
\item \textbf{URL:} \texttt{/category}

\item \textbf{Method:} \texttt{GET}

\item \textbf{URL Parameters:} \texttt{None}

\item \textbf{Query Parameters:} \texttt{None}

\item \textbf{Headers:} \texttt{Accept: application/json}

\item \textbf{Cookies:} \texttt{None}

\item \textbf{Data Parameters:} \texttt{None}

\item \textbf{Success Response:} \texttt{200 OK}

Content:

\begin{lstlisting}
{
  "resource-list": [
    { "name": "Antipasti" },
    { "name": "Primi" },
    { "name": "Secondi" }
  ]
}
\end{lstlisting}

\item \textbf{Error Responses:}

\begin{itemize}[nosep]
\item Common header errors: \texttt{E4V1}, \texttt{E4V2}.
\item \texttt{500 Internal Server Error}, \texttt{E5D1} - inner message \texttt{Database error while retrieving categories.}, inner details are the SQL exception message.
\end{itemize}

\end{itemize}

\subsubsection*{Delete Category}

Deletes a category by name.

\begin{itemize}[nosep]
\item \textbf{URL:} \texttt{/category/\{category\_name\}}

\item \textbf{Method:} \texttt{DELETE}

\item \textbf{URL Parameters:} \texttt{category\_name} - category primary key

\item \textbf{Query Parameters:} \texttt{None}

\item \textbf{Headers:} \texttt{Accept: application/json}

\item \textbf{Cookies:} \texttt{auth\_token} HttpOnly cookie containing a valid JWT with role \texttt{ADMIN}

\item \textbf{Data Parameters:} \texttt{None}

The path parameter is matched as a raw route segment. Database matching is exact because the DAO executes \texttt{DELETE FROM categories WHERE name = -}.

\item \textbf{Success Response:} \texttt{200 OK}

Content:

\begin{lstlisting}
{
  "name": "Pizze"
}
\end{lstlisting}

\item \textbf{Error Responses:}

\begin{itemize}[nosep]
\item Common auth errors: \texttt{401 Unauthorized}, \texttt{E4A2}; \texttt{403 Forbidden}, \texttt{E4A4}.
\item Common header errors: \texttt{E4V1}, \texttt{E4V2}.
\item \texttt{404 Not Found}, \texttt{E4R2} - inner message \texttt{Category 'Pizze' not found.}
\item \texttt{409 Conflict}, \texttt{E5D1} - inner message \texttt{Cannot delete category: it is referenced by existing dishes.}, inner details are the SQL exception message. The schema conflict is detected via SQL state \texttt{23503}, but the handler still emits \texttt{E5D1}.
\item \texttt{500 Internal Server Error}, \texttt{E5D1} - inner message \texttt{Database error while deleting category.}, inner details are the SQL exception message.
\end{itemize}

\end{itemize}

\subsubsection*{Create Ingredient}

Creates an ingredient.

\begin{itemize}[nosep]
\item \textbf{URL:} \texttt{/ingredient}

\item \textbf{Method:} \texttt{POST}

\item \textbf{URL Parameters:} \texttt{None}

\item \textbf{Query Parameters:} \texttt{None}

\item \textbf{Headers:} \texttt{Content-Type: application/json}; \texttt{Accept: application/json}

\item \textbf{Cookies:} \texttt{auth\_token} HttpOnly cookie containing a valid JWT with role \texttt{STAFF} or \texttt{ADMIN}

\item \textbf{Data Parameters:}

\begin{lstlisting}
{
  "name": "Origano",
  "allergens": [],
  "is_frozen": false
}
\end{lstlisting}

The parser accepts \texttt{id}, \texttt{name}, \texttt{allergens}, and \texttt{is\_frozen}. The handler requires only \texttt{name}, checked for \texttt{null} (not blank). \texttt{is\_frozen} and \texttt{allergens} are optional; \texttt{allergens} may be omitted, set to \texttt{null}, or be an empty array. Valid allergen values are the lowercase enum literals \texttt{cereals}, \texttt{crustaceans}, \texttt{eggs}, \texttt{fish}, \texttt{peanuts}, \texttt{soybeans}, \texttt{milk}, \texttt{nuts}, \texttt{celery}, \texttt{mustard}, \texttt{sesame\_seeds}, \texttt{sulphur\_dioxide\_and\_sulphites}, \texttt{lupin}, and \texttt{molluscs}. Invalid allergen enum values are not explicitly handled and fall through to \texttt{E5G1}.

\item \textbf{Success Response:} \texttt{201 Created}

Content:

\begin{lstlisting}
{
  "id": 6,
  "name": "Origano",
  "allergens": [],
  "is_frozen": false
}
\end{lstlisting}

\item \textbf{Error Responses:}

\begin{itemize}[nosep]
\item Common auth errors: \texttt{401 Unauthorized}, \texttt{E4A2}; \texttt{403 Forbidden}, \texttt{E4A4}.
\item Common header errors: \texttt{E4V1}, \texttt{E4V2}, \texttt{E4V3}, \texttt{E4V4}.
\item \texttt{400 Bad Request}, \texttt{E4V6} - inner message \texttt{Missing required fields: name and isFrozen must be provided} when \texttt{name} is \texttt{null}.
\item \texttt{400 Bad Request}, \texttt{E4V8} - inner message \texttt{Malformed JSON in request body.}, inner details are \texttt{Unexpected field: <field>}.
\item \texttt{500 Internal Server Error}, \texttt{E5D1} - inner message \texttt{Failed to create ingredient.} if the DAO returns no row.
\item \texttt{500 Internal Server Error}, \texttt{E5D1} - inner message \texttt{Database error while creating ingredient.}, inner details are the SQL exception message.
\item \texttt{500 Internal Server Error}, \texttt{E5G1} - malformed JSON syntax or invalid allergen enum values fall through to \texttt{Unable to serve the REST request: CREATE\_INGREDIENT.}
\end{itemize}

\item \textbf{Sample Request Body:}

\begin{lstlisting}
{
  "name": "Origano",
  "allergens": [],
  "is_frozen": false
}
\end{lstlisting}
\end{itemize}

\subsubsection*{Get All Ingredients}

Returns all ingredients.

\begin{itemize}[nosep]
\item \textbf{URL:} \texttt{/ingredient}

\item \textbf{Method:} \texttt{GET}

\item \textbf{URL Parameters:} \texttt{None}

\item \textbf{Query Parameters:} \texttt{None}

\item \textbf{Headers:} \texttt{Accept: application/json}

\item \textbf{Cookies:} \texttt{None}

\item \textbf{Data Parameters:} \texttt{None}

The read DAOs use \texttt{ResultSet\#getBoolean("is\_frozen")} without \texttt{wasNull()}, so a database \texttt{NULL} in \texttt{ingredients.is\_frozen} would be serialized as \texttt{false}.

\item \textbf{Success Response:} \texttt{200 OK}

Content:

\begin{lstlisting}
{
  "resource-list": [
    {
      "id": 1,
      "name": "Pomodoro San Marzano",
      "allergens": [],
      "is_frozen": false
    },
    {
      "id": 2,
      "name": "Mozzarella di Bufala",
      "allergens": ["milk"],
      "is_frozen": false
    }
  ]
}
\end{lstlisting}

\item \textbf{Error Responses:}

\begin{itemize}[nosep]
\item Common header errors: \texttt{E4V1}, \texttt{E4V2}.
\item \texttt{500 Internal Server Error}, \texttt{E5D1} - inner message \texttt{Database error while retrieving ingredients}, with no explicit \texttt{error\_details}.
\end{itemize}

\end{itemize}

\subsubsection*{Get Ingredient}

Returns one ingredient by id.

\begin{itemize}[nosep]
\item \textbf{URL:} \texttt{/ingredient/\{ingredient\_id\}}

\item \textbf{Method:} \texttt{GET}

\item \textbf{URL Parameters:} \texttt{ingredient\_id} - integer ingredient id

\item \textbf{Query Parameters:} \texttt{None}

\item \textbf{Headers:} \texttt{Accept: application/json}

\item \textbf{Cookies:} \texttt{None}

\item \textbf{Data Parameters:} \texttt{None}

\item \textbf{Success Response:} \texttt{200 OK}

Content:

\begin{lstlisting}
{
  "id": 2,
  "name": "Mozzarella di Bufala",
  "allergens": ["milk"],
  "is_frozen": false
}
\end{lstlisting}

\item \textbf{Error Responses:}

\begin{itemize}[nosep]
\item Common header errors: \texttt{E4V1}, \texttt{E4V2}.
\item \texttt{400 Bad Request}, \texttt{E4V6} - inner message \texttt{Ingredient id must be a valid integer}
\item \texttt{404 Not Found}, \texttt{E4R2} - inner message \texttt{Ingredient 2 not found in database}
\item \texttt{500 Internal Server Error}, \texttt{E5D1} - inner message \texttt{Database error while retrieving ingredient}, with no explicit \texttt{error\_details}.
\end{itemize}

\end{itemize}

\subsubsection*{Edit Ingredient}

Updates an ingredient.

\begin{itemize}[nosep]
\item \textbf{URL:} \texttt{/ingredient}

\item \textbf{Method:} \texttt{PUT}

\item \textbf{URL Parameters:} \texttt{None}

\item \textbf{Query Parameters:} \texttt{None}

\item \textbf{Headers:} \texttt{Content-Type: application/json}; \texttt{Accept: application/json}

\item \textbf{Cookies:} \texttt{auth\_token} HttpOnly cookie containing a valid JWT with role \texttt{STAFF} or \texttt{ADMIN}

\item \textbf{Data Parameters:}

\begin{lstlisting}
{
  "id": 2,
  "name": "Mozzarella di Bufala DOP",
  "allergens": ["milk"],
  "is_frozen": false
}
\end{lstlisting}

The handler requires \texttt{id}, \texttt{name}, and \texttt{is\_frozen}; unlike create, it also rejects blank \texttt{name}. \texttt{allergens} remains optional. Invalid allergen enum values and malformed JSON syntax are not explicitly mapped and fall through to \texttt{E5G1}.

\item \textbf{Success Response:} \texttt{200 OK}

Content:

\begin{lstlisting}
{
  "id": 2,
  "name": "Mozzarella di Bufala DOP",
  "allergens": ["milk"],
  "is_frozen": false
}
\end{lstlisting}

\item \textbf{Error Responses:}

\begin{itemize}[nosep]
\item Common auth errors: \texttt{401 Unauthorized}, \texttt{E4A2}; \texttt{403 Forbidden}, \texttt{E4A4}.
\item Common header errors: \texttt{E4V1}, \texttt{E4V2}, \texttt{E4V3}, \texttt{E4V4}.
\item \texttt{400 Bad Request}, \texttt{E4V6} - inner message \texttt{Missing required fields: name and isFrozen must be provided}
\item \texttt{400 Bad Request}, \texttt{E4V8} - inner message \texttt{Unexpected key in JSON: Unexpected field: <field>}
\item \texttt{404 Not Found}, \texttt{E4R2} - inner message \texttt{Ingredient with id 2 not found.}
\item \texttt{500 Internal Server Error}, \texttt{E5D1} - inner message \texttt{Database error while editing ingredient.}, inner details are the SQL exception message.
\item \texttt{500 Internal Server Error}, \texttt{E5G1} - malformed JSON syntax or invalid allergen enum values fall through to \texttt{Unable to serve the REST request: EDIT\_INGREDIENT.}
\end{itemize}

\item \textbf{Sample Request Body:}

\begin{lstlisting}
{
  "id": 2,
  "name": "Mozzarella di Bufala DOP",
  "allergens": ["milk"],
  "is_frozen": false
}
\end{lstlisting}
\end{itemize}

\subsubsection*{Delete Ingredient}

Deletes an ingredient by id.

\begin{itemize}[nosep]
\item \textbf{URL:} \texttt{/ingredient/\{ingredient\_id\}}

\item \textbf{Method:} \texttt{DELETE}

\item \textbf{URL Parameters:} \texttt{ingredient\_id} - integer ingredient id

\item \textbf{Query Parameters:} \texttt{None}

\item \textbf{Headers:} \texttt{Accept: application/json}

\item \textbf{Cookies:} \texttt{auth\_token} HttpOnly cookie containing a valid JWT with role \texttt{STAFF} or \texttt{ADMIN}

\item \textbf{Data Parameters:} \texttt{None}

The schema defines \texttt{dish\_ingredients.ingredient\_id} with \texttt{ON DELETE CASCADE}, so deleting an ingredient also deletes its dish bindings.

\item \textbf{Success Response:} \texttt{200 OK}

Content:

\begin{lstlisting}
{
  "id": 6,
  "name": "Origano",
  "allergens": [],
  "is_frozen": false
}
\end{lstlisting}

\item \textbf{Error Responses:}

\begin{itemize}[nosep]
\item Common auth errors: \texttt{401 Unauthorized}, \texttt{E4A2}; \texttt{403 Forbidden}, \texttt{E4A4}.
\item Common header errors: \texttt{E4V1}, \texttt{E4V2}.
\item \texttt{400 Bad Request}, \texttt{E4V6} - inner message \texttt{Invalid ingredient id format: /6} or the actual \texttt{req.getPathInfo()} value, inner details are the \texttt{NumberFormatException} message.
\item \texttt{404 Not Found}, \texttt{E4R2} - inner message \texttt{Ingredient with id 6 not found}
\item \texttt{500 Internal Server Error}, \texttt{E5D1} - inner message \texttt{Unexpected database error: no. <sqlErrorCode>}, inner details are the SQL exception message.
\end{itemize}

\end{itemize}

\subsubsection*{Create Dish}

Creates a dish and its ingredient bindings.

\begin{itemize}[nosep]
\item \textbf{URL:} \texttt{/dish}

\item \textbf{Method:} \texttt{POST}

\item \textbf{URL Parameters:} \texttt{None}

\item \textbf{Query Parameters:} \texttt{None}

\item \textbf{Headers:} \texttt{Content-Type: application/json}; \texttt{Accept: application/json}

\item \textbf{Cookies:} \texttt{auth\_token} HttpOnly cookie containing a valid JWT with role \texttt{STAFF} or \texttt{ADMIN}

\item \textbf{Data Parameters:}

\begin{lstlisting}
{
  "name": "Margherita",
  "description": "Pizza con pomodoro, mozzarella e basilico",
  "price": 8.5,
  "ingredient_ids": [1, 2, 3],
  "category": "Pizze"
}
\end{lstlisting}

Required fields are \texttt{name}, \texttt{price}, \texttt{ingredient\_ids}, and \texttt{category}. \texttt{name} must be non-blank, \texttt{price} must be strictly greater than \texttt{0}, and \texttt{ingredient\_ids} must be a non-empty array. \texttt{description} is optional. The parser also accepts \texttt{id}, but the create logic ignores it. On success, the response contains nested \texttt{ingredients}, not \texttt{ingredient\_ids}. Invalid category names, invalid ingredient ids, and duplicate ingredient ids in the submitted array surface as database errors (\texttt{E5D1}), not as explicit 4xx validation errors.

\item \textbf{Success Response:} \texttt{201 Created}

Content:

\begin{lstlisting}
{
  "id": 6,
  "name": "Margherita",
  "description": "Pizza con pomodoro, mozzarella e basilico",
  "price": 8.5,
  "ingredients": [
    {
      "id": 1,
      "name": "Pomodoro San Marzano",
      "allergens": [],
      "is_frozen": false
    },
    {
      "id": 2,
      "name": "Mozzarella di Bufala",
      "allergens": ["milk"],
      "is_frozen": false
    },
    {
      "id": 3,
      "name": "Basilico fresco",
      "allergens": [],
      "is_frozen": false
    }
  ],
  "category": "Pizze"
}
\end{lstlisting}

\item \textbf{Error Responses:}

\begin{itemize}[nosep]
\item Common auth errors: \texttt{401 Unauthorized}, \texttt{E4A2}; \texttt{403 Forbidden}, \texttt{E4A4}.
\item Common header errors: \texttt{E4V1}, \texttt{E4V2}, \texttt{E4V3}, \texttt{E4V4}.
\item \texttt{400 Bad Request}, \texttt{E4V6} - inner message \texttt{Missing or invalid fields}
\item \texttt{400 Bad Request}, \texttt{E4V8} - inner message \texttt{Malformed JSON in request body.}, inner details are \texttt{Unexpected field: <field>}.
\item \texttt{500 Internal Server Error}, \texttt{E5D1} - inner message \texttt{Failed to create dish.} if the DAO returns no row.
\item \texttt{500 Internal Server Error}, \texttt{E5D1} - inner message \texttt{Database error while creating dish.}, inner details are the SQL exception message.
\item \texttt{500 Internal Server Error}, \texttt{E5G1} - malformed JSON syntax falls through to \texttt{Unable to serve the REST request: CREATE\_DISH.}
\end{itemize}

\item \textbf{Sample Request Body:}

\begin{lstlisting}
{
  "name": "Margherita",
  "description": "Pizza con pomodoro, mozzarella e basilico",
  "price": 8.5,
  "ingredient_ids": [1, 2, 3],
  "category": "Pizze"
}
\end{lstlisting}
\end{itemize}

\subsubsection*{Get All Dishes}

Returns all dishes with nested ingredient objects.

\begin{itemize}[nosep]
\item \textbf{URL:} \texttt{/dish}

\item \textbf{Method:} \texttt{GET}

\item \textbf{URL Parameters:} \texttt{None}

\item \textbf{Query Parameters:} \texttt{None}

\item \textbf{Headers:} \texttt{Accept: application/json}

\item \textbf{Cookies:} \texttt{None}

\item \textbf{Data Parameters:} \texttt{None}

\item \textbf{Success Response:} \texttt{200 OK}

Content:

\begin{lstlisting}
{
  "resource-list": [
    {
      "id": 1,
      "name": "Bruschetta al pomodoro",
      "description": "Pane tostato con pomodoro fresco e basilico",
      "price": 5.5,
      "ingredients": [
        {
          "id": 1,
          "name": "Pomodoro San Marzano",
          "allergens": [],
          "is_frozen": false
        },
        {
          "id": 3,
          "name": "Basilico fresco",
          "allergens": [],
          "is_frozen": false
        }
      ],
      "category": "Antipasti"
    }
  ]
}
\end{lstlisting}

\item \textbf{Error Responses:}

\begin{itemize}[nosep]
\item Common header errors: \texttt{E4V1}, \texttt{E4V2}.
\item \texttt{500 Internal Server Error}, \texttt{E5D1} - inner message \texttt{Database error while retrieving dishes.}, inner details are the SQL exception message.
\end{itemize}

\end{itemize}

\subsubsection*{Get Dish}

Returns one dish by id.

\begin{itemize}[nosep]
\item \textbf{URL:} \texttt{/dish/\{dish\_id\}}

\item \textbf{Method:} \texttt{GET}

\item \textbf{URL Parameters:} \texttt{dish\_id} - integer dish id

\item \textbf{Query Parameters:} \texttt{None}

\item \textbf{Headers:} \texttt{Accept: application/json}

\item \textbf{Cookies:} \texttt{None}

\item \textbf{Data Parameters:} \texttt{None}

\item \textbf{Success Response:} \texttt{200 OK}

Content:

\begin{lstlisting}
{
  "id": 2,
  "name": "Spaghetti al pomodoro",
  "description": "Spaghetti con salsa di pomodoro fresco e basilico",
  "price": 9.0,
  "ingredients": [
    {
      "id": 4,
      "name": "Spaghetti di Gragnano",
      "allergens": ["cereals", "eggs"],
      "is_frozen": false
    },
    {
      "id": 1,
      "name": "Pomodoro San Marzano",
      "allergens": [],
      "is_frozen": false
    }
  ],
  "category": "Primi"
}
\end{lstlisting}

\item \textbf{Error Responses:}

\begin{itemize}[nosep]
\item Common header errors: \texttt{E4V1}, \texttt{E4V2}.
\item \texttt{400 Bad Request}, \texttt{E4V6} - inner message \texttt{Invalid dish id format}, inner details are the \texttt{NumberFormatException} message.
\item \texttt{404 Not Found}, \texttt{E4R2} - inner message \texttt{Dish with id 2 not found.}
\item \texttt{500 Internal Server Error}, \texttt{E5D1} - inner message \texttt{Database error while retrieving dish.}, inner details are the SQL exception message.
\end{itemize}

\end{itemize}

\subsubsection*{Edit Dish}

Updates a dish and replaces its ingredient bindings.

\begin{itemize}[nosep]
\item \textbf{URL:} \texttt{/dish}

\item \textbf{Method:} \texttt{PUT}

\item \textbf{URL Parameters:} \texttt{None}

\item \textbf{Query Parameters:} \texttt{None}

\item \textbf{Headers:} \texttt{Content-Type: application/json}; \texttt{Accept: application/json}

\item \textbf{Cookies:} \texttt{auth\_token} HttpOnly cookie containing a valid JWT with role \texttt{STAFF} or \texttt{ADMIN}

\item \textbf{Data Parameters:}

\begin{lstlisting}
{
  "id": 2,
  "name": "Spaghetti al pomodoro fresco",
  "description": "Spaghetti con salsa di pomodoro e basilico",
  "price": 9.5,
  "ingredient_ids": [4, 1, 3],
  "category": "Primi"
}
\end{lstlisting}

Required fields are \texttt{id}, \texttt{name}, \texttt{price}, \texttt{ingredient\_ids}, and \texttt{category}. \texttt{name} must be non-blank. Unlike create, this handler does not enforce \texttt{price > 0}; it only requires \texttt{price != null}. \texttt{description} is optional. Invalid category names, invalid ingredient ids, and duplicate ingredient ids surface as \texttt{E5D1}. Malformed JSON syntax is not explicitly handled and falls through to \texttt{E5G1}.

\item \textbf{Success Response:} \texttt{200 OK}

Content:

\begin{lstlisting}
{
  "id": 2,
  "name": "Spaghetti al pomodoro fresco",
  "description": "Spaghetti con salsa di pomodoro e basilico",
  "price": 9.5,
  "ingredients": [
    {
      "id": 4,
      "name": "Spaghetti di Gragnano",
      "allergens": ["cereals", "eggs"],
      "is_frozen": false
    },
    {
      "id": 1,
      "name": "Pomodoro San Marzano",
      "allergens": [],
      "is_frozen": false
    },
    {
      "id": 3,
      "name": "Basilico fresco",
      "allergens": [],
      "is_frozen": false
    }
  ],
  "category": "Primi"
}
\end{lstlisting}

\item \textbf{Error Responses:}

\begin{itemize}[nosep]
\item Common auth errors: \texttt{401 Unauthorized}, \texttt{E4A2}; \texttt{403 Forbidden}, \texttt{E4A4}.
\item Common header errors: \texttt{E4V1}, \texttt{E4V2}, \texttt{E4V3}, \texttt{E4V4}.
\item \texttt{400 Bad Request}, \texttt{E4V6} - inner message \texttt{Missing required fields}
\item \texttt{400 Bad Request}, \texttt{E4V8} - inner message \texttt{Malformed JSON in request body.}, inner details are \texttt{Unexpected field: <field>}.
\item \texttt{404 Not Found}, \texttt{E4R2} - inner message \texttt{Dish with id 2 not found.}
\item \texttt{500 Internal Server Error}, \texttt{E5D1} - inner message \texttt{Database error while editing dish.}, inner details are the SQL exception message.
\item \texttt{500 Internal Server Error}, \texttt{E5G1} - malformed JSON syntax falls through to \texttt{Unable to serve the REST request: EDIT\_DISH.}
\end{itemize}

\item \textbf{Sample Request Body:}

\begin{lstlisting}
{
  "id": 2,
  "name": "Spaghetti al pomodoro fresco",
  "description": "Spaghetti con salsa di pomodoro e basilico",
  "price": 9.5,
  "ingredient_ids": [4, 1, 3],
  "category": "Primi"
}
\end{lstlisting}
\end{itemize}

\subsubsection*{Delete Dish}

Deletes a dish by id.

\begin{itemize}[nosep]
\item \textbf{URL:} \texttt{/dish/\{dish\_id\}}

\item \textbf{Method:} \texttt{DELETE}

\item \textbf{URL Parameters:} \texttt{dish\_id} - integer dish id

\item \textbf{Query Parameters:} \texttt{None}

\item \textbf{Headers:} \texttt{Accept: application/json}

\item \textbf{Cookies:} \texttt{auth\_token} HttpOnly cookie containing a valid JWT with role \texttt{STAFF} or \texttt{ADMIN}

\item \textbf{Data Parameters:} \texttt{None}

The schema keeps \texttt{order\_dishes.dish\_id} as a foreign key without \texttt{ON DELETE CASCADE}, so deleting a dish referenced by an order fails with a database error (\texttt{E5D1}), not a conflict code.

\item \textbf{Success Response:} \texttt{204 No Content}

Content: \texttt{None}

\item \textbf{Error Responses:}

\begin{itemize}[nosep]
\item Common auth errors: \texttt{401 Unauthorized}, \texttt{E4A2}; \texttt{403 Forbidden}, \texttt{E4A4}.
\item Common header errors: \texttt{E4V1}, \texttt{E4V2}.
\item \texttt{400 Bad Request}, \texttt{E4V6} - inner message \texttt{Invalid dish id format}, inner details are the \texttt{NumberFormatException} message.
\item \texttt{404 Not Found}, \texttt{E4R2} - inner message \texttt{Dish with id 2 not found.}
\item \texttt{500 Internal Server Error}, \texttt{E5D1} - inner message \texttt{Database error while deleting dish.}, inner details are the SQL exception message.
\end{itemize}

\end{itemize}
